% Part: lambda-calculus
% Chapter: introduction
% Section: overview

\documentclass[../../../include/open-logic-section]{subfiles}

\begin{document}

\olfileid[id]{lam}{int}{ovr}
\olsection{Ikhtisar}

Kalkulus lambda mula-mula dirancang oleh Alonzo Church pada awal 1930-an
sebagai landasan bagi logika konstruktif, dan \emph{bukan} sebagai model
fungsi-fungsi yang dapat dihitung. Namun, kalkulus ini segera terbukti
ekuivalen dengan definisi komputabilitas lainnya, seperti fungsi-fungsi yang
dapat dihitung oleh Turing dan fungsi-fungsi rekursif parsial. Kenyataan bahwa
hasil ini pada mulanya agak mengejutkan justru membuat karakterisasi tersebut
semakin menarik.

Notasi lambda merupakan cara praktis untuk merujuk langsung kepada suatu
fungsi melalui ekspresi simbolik yang mendefinisikannya, alih-alih memberi
nama kepada fungsi itu terlebih dahulu. Daripada mengatakan ``misalkan $f$
adalah fungsi yang didefinisikan oleh $f(x)=x+3$,'' kita dapat mengatakan
``misalkan $f$ adalah fungsi $\lambd[x][(x+3)]$.'' Dengan kata lain,
$\lambd[x][(x+3)]$ hanyalah sebuah \emph{nama} bagi fungsi yang menambahkan
tiga pada argumennya. Dalam ekspresi ini, $x$ adalah variabel semu atau
penanda tempat: fungsi yang sama dapat pula dinyatakan dengan
$\lambd[y][(y+3)]$. Notasi ini tetap berfungsi apabila terdapat parameter
lain. Misalkan, sebagai contoh, $g(x,y)$ adalah fungsi dua variabel dan $k$
suatu bilangan asli. Maka $\lambd[x][g(x,k)]$ adalah fungsi yang memetakan
setiap~$x$ ke~$g(x,k)$.

Cara mendefinisikan fungsi dari ekspresi simbolik ini dikenal sebagai
\emph{abstraksi lambda}. Sisi lain abstraksi lambda adalah \emph{aplikasi}:
andaikan kita mempunyai suatu fungsi~$f$ (misalnya, yang didefinisikan pada
bilangan asli), kita dapat \emph{menerapkan} fungsi itu pada nilai apa pun,
misalnya~2. Dalam notasi biasa, tentu saja, kita menulis~$f(2)$ untuk hasilnya.

Apa yang terjadi apabila abstraksi lambda digabungkan dengan aplikasi?
Ekspresi yang dihasilkan dapat disederhanakan dengan ``memasukkan'' argumen
yang diterapkan sebagai pengganti variabel yang diabstraksikan. Sebagai contoh,
\[
(\lambd[x][(x + 3)])(2)
\]
dapat disederhanakan menjadi~$2+3$.

Sampai tahap ini, kita baru memperkenalkan notasi baru bagi gagasan-gagasan
biasa. Akan tetapi, kalkulus lambda menyimpang secara lebih radikal dari sudut
pandang teori himpunan. Dalam kerangka ini:
\begin{enumerate}
\item Segala sesuatu menyatakan suatu fungsi.
\item Fungsi dapat didefinisikan dengan abstraksi lambda.
\item Apa pun dapat diterapkan pada apa pun yang lain.
\end{enumerate}
Sebagai contoh, jika $F$ adalah suatu suku dalam kalkulus lambda, $F(F)$ selalu
dianggap bermakna. Kerangka yang bebas ini dikenal sebagai kalkulus lambda
\emph{tanpa tipe}; ``tanpa tipe'' berarti ``tidak ada pembatasan mengenai apa
yang dapat diterapkan pada apa.''

\begin{digress}
Ada pula kalkulus lambda \emph{bertipe}, suatu variasi penting dari versi
tanpa tipe. Meskipun dalam banyak hal kalkulus lambda bertipe serupa dengan
versi tanpa tipe, kalkulus ini jauh lebih mudah diselaraskan dengan kerangka
teori himpunan klasik dan memiliki sejumlah sifat yang sangat berbeda.

Penelitian mengenai kalkulus lambda terbukti memainkan peran sentral dalam
ilmu komputer teoretis dan dalam perancangan bahasa pemrograman. LISP, yang
dirancang oleh John McCarthy pada 1950-an, merupakan contoh awal bahasa yang
dipengaruhi oleh gagasan-gagasan ini.
\end{digress}

\end{document}
